\subsection{Linear effects: importance is in squared-coefficient units}
\label{sec:linear}


Suppose that a random forest is fitted to the data: 
$y = a + 2b + 3c + \dots + 11k$ (11 features, coefficients $1\ldots11$)
$a,\dots,k \sim U(0,1)$, $n = 1000$.
Over seeds $1$--$10$, we regressing importance on $\beta^2$
gives a mean
$R^2$ of $0.98$ (range $0.96$--$0.99$) against $0.91$ on $\beta$ (range
$0.88$--$0.94$), and $\beta^2$ is the better fit on every seed;
Figure~\ref{fig:superlinear} shows one of them. 


\begin{figure}[htb]
\centering
\includegraphics[width=0.95\textwidth]{fig1_superlinear.png}
\caption{Impurity importance against the coefficient $\beta$ (left, curved) and
against the variance contribution $\beta^2$ (right, linear, $R^2=0.99$).
Importance is linear in $\beta^2$, not in $\beta$.}
\label{fig:superlinear}
\end{figure}

%This empirically confirms that MDI importance is denominated in
%squared-coefficient (variance) units rather than coefficient units,
%as predicted by Scornet's identity.

For a linear term the contributed
variance is $\beta_j^2\,\mathrm{Var}(X_j)$, and here the standard-uniform
variance $\mathrm{Var}(X_j)=1/12$ is
common to all features, so the MDI should track $\beta_j^2$.

The experiment supports that, but the fit is not \emph{perfect}. This is also 
observed from the ranking of the feature importances: MDI only ranks
feature constributions correctly when they differ enough and there is enough data
to dwarf the randommness of the procedure of fitting the trees.
For this data, 
$\beta = 1\ldots11$ at the default \code{mtry}, only $3$ of $10$ seeds rank
all eleven correctly (mean Spearman $0.99$). 
AIR spreads weak features (those with smaller coefficients) further apart,
% (the $\beta = 5$ feature scores $16.4$ times the $\beta = 1$ feature, against
% $2.45$ under plain impurity, where the variance identity predicts $25$),
but it is also
far noisier from seed to seed, and the two effects offset: it ranks no better than MDI on these data
(mean Spearman $0.99$). Even where
the order holds, ``how much more'' is not comparable to a coefficient.

\subsubsection{Relative interpretations of feature importances}
An experiment that examines divergences from the expected ratios between 
variable importances enables determining in more detail how and why 
variable importances in practice differ from the theoretical expectation
of them tracking  $\beta_j^2$.
Here, 
a random forest is fitted a data where feature contributions are doubled per term: $y = a + 2b + 4c$, with
$a,b,c \sim U(0,1)$, $n = 1000$. The forest's out-of-bag (OOB) $R^2$ is $0.98$.
Impurity importance (both MDI and AIR) ranks the features in the correct
order, 
but their \emph{ratios} do not match the ordinary-least-squares (OLS)
coefficient ratios (table entries are means over seeds $1$--$10$, ranges in
parentheses):

\begin{center}
\begin{tabular}{rccc}
\toprule
feature & OLS coefficient & MDI (ratio) & AIR (ratio)
\\
\midrule
a & 1 & 1.00 & 1.00 \\
b & 2 & 2.24 (2.00--2.47) & 4.19 (2.97--5.53) \\
c & 4 & 7.19 (6.49--8.21) & 16.88 (11.06--22.94) \\
\bottomrule
\end{tabular}
\end{center}

A coefficient ratio of $1:2:4$ becomes an importance ratio of $1:2.2:7.2$ under
MDI and $1:4.2:16.9$ under AIR: the
feature ordering is correct and the growth is clearly super-linear, but the MDI spread is not
 $1:4:16$, the ratio of the \emph{squared} coefficients that Scornet's identity in 
Section~\ref{sec:background-strands} predicts. 

%For a linear term the contributed
%variance is $\beta_j^2\,\mathrm{Var}(X_j)$, and here the standard-uniform
%variance $\mathrm{Var}(X_j)=1/12$ is
%common to all features, so the MDI should track $\beta_j^2$.k

The difference 
%to this expectation
is due to the number of features drawn as
candidates at each node (a configurable parameter called \code{mtry}), intended
to decorrelate the trees. In \code{ranger} it defaults
to $\lfloor \sqrt{p} \rfloor$, here $1$: each node must split on the
single feature it drew, so $a$ takes splits it would never win in competition
and
is inflated at the expense of $c$. Letting features compete removes the dilution
(means over seeds $1$--$10$, $n=1000$, 500 trees):

\begin{center}
\begin{tabular}{clc}
\toprule
\code{mtry} & MDI (ratio) & \\
\midrule
1 (default at $p=3$) & $1:2.25:7.26$ & compressed \\
2                    & $1:3.93:15.5$ & $\approx$ predicted $1:4:16$ \\
3 ($=p$, no split randomisation) & $1:4.45:19.1$ & overshoots \\
\bottomrule
\end{tabular}
\end{center}

At $\code{mtry}=p$ the
strongest feature wins the root split in every tree, and overshoots. Tree depth
matters too, but not in the direction Scornet's truncated-tree condition might
suggest. He proves the identity for truncated trees and warns that fully grown
trees yield positively biased MDI \citep{scornet2023}, and explains that the
bias comes from fitting the noise term; $y$ here is noiseless, so it cannot
operate. A sweep limiting the maximal depth of trees
%\code{max.depth} sweep (seeds $1$--$10$) 
shows the opposite, finite-sample
effect: truncating the trees removes splits, and the features with the smallest
variance contributions are the first to lose them. At depth $2$ with
$\code{mtry}=p$
the weakest feature earns no splits at all. Deepening lowers the ratios in
both settings: at $\code{mtry}=p$ they approach the predicted $1:4:16$
($1:12.0:58.3$ at depth $4$, $1:4.45:19.1$ fully grown), and at the default
$\code{mtry}=1$ they fall further below it ($1:3.38:12.1$ at depth $2$,
$1:2.25:7.26$ fully grown). The residual overshoot at $\code{mtry}=p$ is a
finite-sample effect: it diminishes as $n$ grows
($1:4.15:17.1$ at $n=4000$, $1:4.01:16.1$ at $n=64000$; means over seeds
$1$--$10$). This is aligned with Scornet's theorem about greedy CART splitting
being asymptotic in $n$. The compression at $\code{mtry}=1$, which the theorem
does not
cover, also shrinks with $n$ but far more slowly ($1:3.51:13.7$ at
$n=64000$). The deviations in the table are finite-sample properties of the
fitting procedure --- split randomisation and greedy split selection --- and
shrink toward the squared-coefficient prediction as $n$ grows.

%The AIR bias correction behaves differently. 
%It
%pairs each feature with a reordered copy of itself --- same distribution,
%same number of split points, no relationship to the response --- and
%subtracts the variance reduction that copy earns from the feature's own.
%It is therefore a difference of variance reductions, expressed in units of
%variance, and no closer to a per-unit effect size than plain impurity.
%
Compared to MDI, AIR is nearly
\code{mtry}-insensitive: averaged over 10 randomly seeded experiment runs, its ratios are
$1:4.2:16.9$ at the
default, $1:4.0:16.0$ at $\code{mtry}=2$, and $1:4.1:16.7$ at $\code{mtry}=3$,
near the predicted $1:4:16$ throughout.
%
%The price is sampling variance, and it comes from that same subtraction.
Note, the AIR formula includes and additional term (a re-ordered copy of the feature) compared to MDI which 
substracts the reordered variance reduction from the features own.
In this experiment, for the weakest feature the two terms in AIR nearly cancel, leaving a small difference
between two large numbers, so the ordinary noise each carries falls on a much
smaller result: feature $a$'s AIR varies about three times as much from seed to
seed as plain impurity's ($27\%$ against $10\%$ relative standard deviation).
The strongest feature, where the subtraction takes little off the total, stays
steady at $6\%$. Since the unstable feature is the one the ratio divides by,
AIR's $c$ ratio spans $11.1$--$22.9$ over the ten seeds against $6.5$--$8.2$
for plain impurity. The increased sampling variance suggest that AIR ratios 
are readable only as an average over seeds.

%\paragraph{Takeaway.}
%Furthermore, an interaction lets a variable with no marginal effect outrank one
%that has
%(Section~\ref{sec:interactions}), and unequal feature distributions reorder
%the ranking on their own (Section~\ref{sec:confound}).

\paragraph{Practical implications.} Read an impurity ranking as an ordering
among comparably distributed features, and expect inversions where
contributions are close. Its ratios are in squared-coefficient units and
further compressed by the default \code{mtry}, so a bar twice as tall does not
mean an effect twice as large. AIR was built to remove the bias
toward null and high-cardinality features, but for continous features, it also follows the
variance contribution more closely than plain MDI, but varies far more from seed to seed on
weak features; average it over seeds before reading a ratio from it.
